% Pollis Domain Packages: Market Event Studies (PollisEventStudies.jl) -- draft 1 (22 October 2026)
% Compile with: latexmk -pdf dp_event_studies.tex && latexmk -c
% Shared look for the Pollis New Methods documents (same fonts and colours as the brochure and the workflow handouts).
\documentclass[10pt,a4paper]{article}

\usepackage[a4paper,left=18mm,right=18mm,top=15mm,bottom=17mm,footskip=8mm]{geometry}
\usepackage[T1]{fontenc}
\usepackage[default,semibold]{sourcesanspro}
\usepackage{textcomp}
\usepackage[italic]{mathastext}
\usepackage{amsmath}
\usepackage{microtype}
\usepackage[table]{xcolor}
\usepackage{tikz}
\usepackage[most]{tcolorbox}
\usepackage{fontawesome5}
\usepackage{tabularx,booktabs}
\usepackage{enumitem}
\usepackage{fancyhdr}
\usepackage[hidelinks]{hyperref}

% ---------------------------------------------------------------- colours (as in the brochure)
\definecolor{navy}{HTML}{1C3A5E}
\definecolor{pblue}{HTML}{3B78B5}
\definecolor{amber}{HTML}{F2B33D}
\definecolor{deepamber}{HTML}{D9902A}
\definecolor{pale}{HTML}{EEF3F9}
\definecolor{ink}{HTML}{2A2F36}
\definecolor{logonavy}{HTML}{173D63}
\definecolor{logodark}{HTML}{C98A1E}
\definecolor{logoamber}{HTML}{F2B446}
\definecolor{logocream}{HTML}{FEEFD2}
\definecolor{soft}{HTML}{5E6773}
\definecolor{agree}{HTML}{2F7D4F}
\definecolor{differ}{HTML}{B23B3B}
\color{ink}

\setlength{\parindent}{0pt}
\setlength{\parskip}{4pt}
\setlist{nosep,leftmargin=*,itemsep=2pt}
\renewcommand{\arraystretch}{1.15}

\pagestyle{fancy}
\fancyhf{}
\renewcommand{\headrulewidth}{0pt}
\fancyfoot[L]{\footnotesize\color{soft}\textbf{Pollis}\enspace New Methods\enspace·\enspace\nmtheme}
\fancyfoot[R]{\footnotesize\color{soft}\thepage}
\newcommand{\nmtheme}{}
\newcommand{\nmdate}{11 October 2026}

% The Pollis logo, redrawn so it prints sharply at any size
\newcommand{\pollislogo}[1]{%
  \begin{tikzpicture}[x=#1/512,y=-#1/512,baseline=(current bounding box.center)]
    \fill[logonavy] (0,0) rectangle (512,512);
    \fill[logodark] (36,422) -- (104,256) -- (172,422) -- cycle;
    \fill[logoamber] (104,422) -- (203,180) -- (302,422) -- cycle;
    \fill[logocream] (203,422) -- (339,89) -- (476,422) -- cycle;
  \end{tikzpicture}}

% Top of the first page: logo, name; then the kicker, title and a one-line promise
\newcommand{\nmhead}[3]{%
  \renewcommand{\nmtheme}{#1}%
  \thispagestyle{fancy}%
  \begin{minipage}[c]{13mm}\pollislogo{11mm}\end{minipage}%
  \begin{minipage}[c]{\dimexpr\linewidth-13mm\relax}
    {\color{navy}\Large\bfseries Pollis}\enspace{\color{soft}New Methods}\hfill
    {\color{deepamber}\small\bfseries\MakeUppercase{#1}}
  \end{minipage}\par\vspace{7mm}
  {\color{navy}\fontsize{22}{26}\selectfont\bfseries #2}\par\vspace{4pt}
  {\color{soft}\large #3}\par\vspace{5mm}}

\newcommand{\nmsection}[1]{\par\vspace{5pt}{\color{navy}\large\bfseries #1}\par\vspace{1pt}}
\newcommand{\nmsub}[1]{\par\vspace{2pt}{\color{deepamber}\bfseries #1}\par}
\newcommand{\gt}{\,{\color{deepamber}\bfseries ›}\,}
\newcommand{\menupath}[1]{{\bfseries\color{navy}#1}}
\newcommand{\fn}[1]{\texttt{\small #1}}
% Citations: \nmcite{key}{Breiman (2001)} links to the entry \nmref{key}{...} in the References
\newcommand{\nmcite}[2]{\hyperlink{ref:#1}{\color{pblue}#2}}
\newcommand{\nmref}[2]{\raisebox{1.1em}[0pt][0pt]{\hypertarget{ref:#1}{}}#2\par}
\newcommand{\nmdoi}[1]{\href{https://doi.org/#1}{\color{pblue}doi:#1}}
% Where a notebook is in Pollis: its title in Browse Methods, and its file
\newcommand{\browse}[2]{\menupath{Compose \gt New Methods \gt Browse Methods}: \emph{#1} (\fn{#2.ipynb})}

% Key facts at the top: source, missing in Julia, notebook, result
\newcommand{\factsheet}[4]{%
  \begin{tcolorbox}[enhanced, colback=pale, colframe=pale, boxrule=0pt, sharp corners,
    left=7pt, right=7pt, top=5pt, bottom=5pt, fontupper=\small]
  \begin{tabularx}{\linewidth}{@{}>{\leavevmode\bfseries\color{navy}}p{27mm}>{\leavevmode}X@{}}
  Tested against & #1\\
  Missing in Julia & #2\\
  Notebook & #3\\
  Result & #4
  \end{tabularx}
  \end{tcolorbox}}

% Boxes
\newtcolorbox{pbox}[2][]{enhanced, colback=pale, colframe=amber, boxrule=0pt,
  leftrule=3pt, sharp corners, left=7pt, right=7pt, top=6pt, bottom=6pt,
  fontupper=\small\raggedright, before upper={{\color{navy}\normalsize\bfseries #2}\par\smallskip}, #1}
\newtcolorbox{nbox}[2][]{enhanced, colback=navy, colframe=navy, boxrule=0pt,
  sharp corners, left=8pt, right=8pt, top=7pt, bottom=7pt, coltext=white,
  fontupper=\small\raggedright, before upper={{\color{amber}\normalsize\bfseries #2}\par\smallskip}, #1}

% Value tables: published beside Julia
\newenvironment{valuetable}[1]{%
  \par\small\tabularx{\linewidth}{#1}\toprule}{\bottomrule\endtabularx\par}
\newcommand{\ok}{{\color{agree}\faCheck}}
\newcommand{\no}{{\color{differ}\faTimes}}

% End of the document: credit and licence
\newcommand{\nmcredit}{%
  \par\vfill
  \begin{minipage}[b]{0.78\linewidth}\footnotesize\color{soft}
  Methods written and tested with Claude Code (Anthropic) in Pollis, from the published papers and their data.\par
  \textcopyright\ 2026 Antonio Saragga Seabra\enspace·\enspace Licensed under \href{https://creativecommons.org/licenses/by-sa/4.0/}{CC BY-SA 4.0}\enspace·\enspace Draft, \nmdate
  \end{minipage}\hfill
  \begin{minipage}[b]{0.18\linewidth}\raggedleft\pollislogo{9mm}\par\vspace{2pt}{\footnotesize contact@pollis.org}\end{minipage}}


% A package Pollis needs for a feature of its own, not one that reproduces a paper: its own header, key facts and credit
\renewcommand{\nmdate}{22 October 2026}
\fancyfoot[L]{\footnotesize\color{soft}\textbf{Pollis}\enspace Domain Packages\enspace·\enspace\nmtheme}
\renewcommand{\nmhead}[3]{%
  \renewcommand{\nmtheme}{#1}%
  \thispagestyle{fancy}%
  \begin{minipage}[c]{13mm}\pollislogo{11mm}\end{minipage}%
  \begin{minipage}[c]{\dimexpr\linewidth-13mm\relax}
    {\color{navy}\Large\bfseries Pollis}\enspace{\color{soft}Domain Packages}\hfill
    {\color{deepamber}\small\bfseries\MakeUppercase{#1}}
  \end{minipage}\par\vspace{7mm}
  {\color{navy}\fontsize{22}{26}\selectfont\bfseries #2}\par\vspace{4pt}
  {\color{soft}\large #3}\par\vspace{5mm}}
\renewcommand{\factsheet}[4]{%
  \begin{tcolorbox}[enhanced, colback=pale, colframe=pale, boxrule=0pt, sharp corners,
    left=7pt, right=7pt, top=5pt, bottom=5pt, fontupper=\small]
  \begin{tabularx}{\linewidth}{@{}>{\leavevmode\bfseries\color{navy}}p{27mm}>{\leavevmode}X@{}}
  In Pollis & #1\\
  Missing in Julia & #2\\
  Package & #3\\
  Tests & #4
  \end{tabularx}
  \end{tcolorbox}}
\renewcommand{\nmcredit}{%
  \par\vfill
  \begin{minipage}[b]{0.78\linewidth}\footnotesize\color{soft}
  Package written and tested with Claude Code (Anthropic) in Pollis.\par
  \textcopyright\ 2026 Antonio Saragga Seabra\enspace·\enspace Licensed under \href{https://creativecommons.org/licenses/by-sa/4.0/}{CC BY-SA 4.0}\enspace·\enspace Draft, \nmdate
  \end{minipage}\hfill
  \begin{minipage}[b]{0.18\linewidth}\raggedleft\pollislogo{9mm}\par\vspace{2pt}{\footnotesize contact@pollis.org}\end{minipage}}

\begin{document}

\nmhead{Event studies}
  {Event Studies: a filing's minutes, a law's years}
  {How did the market take an 8-K, or the Fed's statement, in the minutes after it came out? And what did a rule, a law
   or a rollout change, period by period, where it came first?}

\factsheet
  {\menupath{Compose \gt Workflows \gt Event Studies}: studies of the filings an alert on SEC filings fires, and of the
   FOMC's releases. Pollis captures the minute prices into \fn{\textasciitilde/.pollis/event-studies}, draws the CAAR as
   they come in, and writes a \fn{study.ipynb} that runs the package.}
  {intraday event studies from minute bars: periods measured in minutes, the next open for a filing outside the session,
   trading that is not every minute, and the moves of rate futures at a release. The Julia event-study code we compared
   works on daily returns. Also the stepped-wedge mixed model with effects by exposure time, which R's steppedwedge
   has.}
  {\fn{PollisEventStudies.jl}, registered in PollisRegistry: \fn{run\_study("path/to/study")}.}
  {116 tests: synthetic studies whose answers are known, built minute by minute, and the panel estimators against R's
   fixest, did, fwildclusterboot and steppedwedge; no published intraday table yet (see Still open).}

\nmsection{Why a package}
The packages in Pollis so far each came from a replication: a paper, its data and its printed results. This one comes
from a feature: Pollis watches EDGAR and the Federal Reserve, and an event study turns an alert's filings into evidence.
The analysis has to live somewhere that a notebook, a test suite and other researchers can all reach.

\nmsub{Where event studies are used}
An event reaches each unit at a date of its own, and the question is what changed after it, period by period.
\begin{itemize}
  \item \textbf{Finance: a company's news, firms under a new rule.} The minutes after a filing or the Fed's statement,
    the subject of the first part of this document; and the years after a rule reaches some firms first, such as the
    business combination laws that states passed one after another from the mid-1980s, shielding the firms
    incorporated in them from takeovers (\nmcite{bm2003}{Bertrand and Mullainathan 2003}).
  \item \textbf{Policy: states passing a law.} The expansion of Medicaid under the Affordable Care Act, adopted by some
    states in 2014 and by others later or not at all, and the mortality of the near-elderly low-income adults it
    covered (\nmcite{mjw2021}{Miller, Johnson and Wherry 2021}).
  \item \textbf{Health: sites joining a programme.} Hospitals or clinics switch to a new treatment or protocol a group at
    a time, in an order drawn at random: the stepped-wedge cluster randomized trial (\nmcite{hh2007}{Hussey and Hughes
    2007}).
  \item \textbf{Education: districts reforming school finance.} Court-ordered reforms changed school spending in
    different states in different years, and the children in school after them were followed into adulthood
    (\nmcite{jjp2016}{Jackson, Johnson and Persico 2016}).
  \item \textbf{Technology: a product reaching one market after another.} Facebook opened to US colleges one after
    another from 2004, and students' mental health worsened after it arrived (\nmcite{blm2022}{Braghieri, Levy and
    Makarin 2022}).
\end{itemize}

\nmsub{Three designs}
What stands in for the outcome without the event tells them apart; the estimators follow from it.
\begin{itemize}
  \item \textbf{A model of the normal outcome, no control units.} The market event study
    (\nmcite{ffjr1969}{Fama, Fisher, Jensen and Roll 1969}; \nmcite{mackinlay1997}{MacKinlay 1997}): each security's
    return around the event against a model of its normal return, estimated before it. It needs an event sharp enough,
    and a window short enough, for nothing else to move the price. Its counterpart in health is the interrupted time
    series, a segmented regression of one series before and after an intervention (\nmcite{lcg2017}{Lopez Bernal,
    Cummins and Gasparrini 2017}), which this package does not have. The first part of this document.
  \item \textbf{Control units, timing not drawn at random.} The panel event study, or dynamic difference-in-differences:
    units treated at different dates against units not yet or never treated, under parallel trends (without the
    treatment, the treated units' outcomes would have moved as the controls' did). The policy, education and technology
    examples above are of this kind, and so are observational health studies of a rollout. Two-way fixed effects with
    an indicator of each event time, and the estimators of \nmcite{sun2021}{Sun and Abraham (2021)} and
    \nmcite{cs2021}{Callaway and Sant'Anna (2021)}: the second part, from Panel event studies on.
  \item \textbf{Control units, timing drawn at random.} The stepped-wedge trial makes the same comparison, but the order
    is drawn at random, so parallel trends hold by design rather than by assumption. Its usual analysis is a mixed model
    with an effect of each period and a random effect of each cluster (\nmcite{hh2007}{Hussey and Hughes 2007}), which
    takes the full effect from the first period of treatment; when the effect grows with the time on treatment, that
    model is biased, as two-way fixed effects are (\nmcite{kvxhh2022}{Kenny et al. 2022}). The estimators of the second
    part remain valid under random timing, and \nmcite{rs2023}{Roth and Sant'Anna (2023)} give the most efficient ones
    under it. The package fits the mixed model too, with Kenny et al.'s effects by exposure time: Stepped wedges, below.
\end{itemize}

\nmsection{The data}
A study is a folder: \fn{study.toml} (the benchmark, \fn{\^{}GSPC} by default; the model; the estimation days, 5 by
default; the periods, 30 minutes before to 120 after and the next open by default; and \fn{group\_by}),
\fn{events.csv} (a filing a row, with its form, 8-K items, EDGAR's acceptance time, the event's time, whether it came
during the regular session and the next open after it) and \fn{prices/} (a file per filing with the minute bars of the
stock and of the benchmark). The event's time is the earlier of the newswire release and EDGAR's acceptance. A minute in
which the stock did not trade has no bar.

\nmsection{The method}
\nmsub{Normal returns}
\emph{Market-adjusted}: the abnormal return is the stock's log return less the benchmark's. \emph{Market model}: alpha
and beta are estimated on the regular sessions of the estimation days before the period. A thinly traded stock is
measured from one trade to the next, against the benchmark's move over the same minutes, and a return over $g$ minutes
is weighed by $1/g$ (weighted least squares; ordinary least squares when the stock trades every minute). A return of
zero in each minute without a trade would bias beta towards zero. This is the treatment of thin trading in Stata's
eventstudy2 and in EventStudy.jl.

\nmsub{Cumulative abnormal returns}
Each event's CAR is followed a minute at a time from the start of the period, from the filing and, for a filing outside
the regular session, from the next open, where the market first trades on it. The AAR and CAAR a minute at a time come
with their cross-sectional t. A filing outside the session, in a period that follows the next open, is tested at the end
of the period's minutes after the next open: its CAR from the filing would end while the market is closed. Events
still capturing prices count up to the last complete minute.

\nmsub{Tests at the end of the period, each with its two-sided p-value}
\begin{itemize}
  \item \textbf{The CAR's cross-sectional t.}
  \item \textbf{Boehmer, Musumeci and Poulsen (1991)}: the t of the standardized CARs,
    $\mathrm{SCAR}_i = \mathrm{CAR}_i / (\hat\sigma_i \sqrt{L_i})$, with $\hat\sigma_i$ the standard deviation of a
    minute's abnormal return on the estimation days and $L_i$ the minutes of the period in which the stock or the
    benchmark has a bar (every minute of the session, whether the stock traded or not). Robust to the volatility an
    event brings.
  \item \textbf{Adjusted for cross-correlation} (\nmcite{kp2010}{Kolari and Pynn\"onen 2010}): events at the same time
    move together, and the BMP t overstates the evidence. It is multiplied by $\sqrt{(1-\bar r)/(1+(n-1)\bar r)}$,
    $\bar r$ the mean correlation over all pairs: that of their standardized abnormal returns over the estimation
    minutes they share when their periods overlap (30 at least), zero when they do not.
  \item \textbf{The generalized rank test} (\nmcite{kp2011}{Kolari and Pynn\"onen 2011}): each event's SCAR, divided by
    the SCARs' cross-sectional standard deviation, is ranked among its standardized abnormal returns on the estimation
    days. Robust as well to returns far from normal; a t with $T-2$ degrees of freedom.
  \item \textbf{Abnormal volume}: the mean $\log(1+\text{volume})$ of the session's minutes after the event (or the next
    open) less that of its minutes before, with its cross-sectional t.
\end{itemize}
The p-values come from Student's t (Distributions.jl), with $n-1$ degrees of freedom for the cross-sectional tests.

\nmsub{Groups}
With \fn{group\_by} (\fn{"form"}, \fn{"item"}, \fn{"session"}, \fn{"ticker"}, or a function of the event) every result
above is given for each group too; a filing with several 8-K items is in the group of each. Pollis sets it from the
study's Group By field. Grouping follows EventStudy.jl.

\nmsub{FOMC releases}
For the statement, projections and minutes, the move of each instrument over each period (10 minutes before to 20 after
by default): federal funds and SOFR futures in basis points of the implied rate, FF1 scaled by $D/(D-d)$, $D$ the days
of the month and $d$ the release's day, and FF2 in its place in the month's last week (\nmcite{kuttner2001}{Kuttner
2001}); the 5- and 10-year yields in basis points; prices in per cent. The mean move with its cross-sectional t, and the
mean path a minute at a time.

\nmsection{A worked example}
\nmsub{The study}
Pollis writes \fn{study.toml} when the study is made on the Event Studies page, with a comment over each setting. A study
of the filings of an alert on 8-Ks, with the market model, grouped by session:

\begingroup\footnotesize
\begin{verbatim}
title = "Earnings 8-Ks"
alert = "earnings-8k"
benchmark = "^GSPC"
model = "marketModel"
event_time = "firstPublic"
estimation_days = 5
group_by = "session"

[[periods]]
name = "Default"
before = 30
after = 120
next_open = true

[[periods]]
name = "Short Run"
before = 5
after = 15
next_open = false
\end{verbatim}
\endgroup


\nmsub{The report}
\fn{run\_study(folder)} reads the folder, analyses each period, prints the report and writes \fn{results/}. Below, the
report of the package's synthetic study of three filings (the first of the tests), of its Default period: one filing
during the session that jumps 2\%, one after Friday's close that falls 3\% at Monday's open, and one during the session
that jumps 1\%. The AAR and CAAR tables are printed at minutes $-30$, $-15$, $-5$, 0, 1, 5, 15, 30, 60 and 120; the
files hold every minute.

\begingroup\footnotesize
\begin{verbatim}
Synthetic: benchmark ^GSPC, market model on the 5 trading days before; 3 events with a ticker
and an acceptance time

Default: 30 minutes before the filing to 120 after; 3 events with prices

Minutes from the filing
  minute  events      AAR %        t     CAAR %        t
     -30       3                          0.000
      -5       3     -0.067   -1.684     -0.023   -1.549
       0       3      0.011    0.644     -0.045   -0.254
       1       3      1.116    1.785      1.071    1.961
      15       3     -0.010   -0.386      0.937    1.986
     120       3      0.086    1.689      1.618    1.997

Minutes from the next open, 1 events
  minute  events      AAR %        t     CAAR %        t
       0       1                          0.000
       1       1     -2.974              -2.974
     120       1     -0.062              -1.919

At +120 minutes: CAR 0.978% (t = 0.675, p = 0.569, 3 events)
Boehmer, Musumeci and Poulsen: mean standardized CAR 0.737 (t = 0.600, p = 0.609, 3 events)
  adjusted for cross-correlation (Kolari and Pynnonen 2010, mean correlation -0.008): t = 0.608,
  p = 0.605
Generalized rank test (Kolari and Pynnonen 2011): t = 0.713, p = 0.476
Abnormal volume: 0.811 (t = 91.512, p = 0.000, 3 events)

By session, at +120 minutes
session                   events      CAR %        t    BMP t     KP t  GRANK t     volume        t
During the session             2      2.426   31.783   34.831   34.831    2.458      0.808   55.958
Outside the session            1     -1.919                                          0.817
\end{verbatim}
\endgroup

The filing after the close is tested on the 120 minutes after Monday's open, where the market first traded on it:
$-1.919$\%, the last row of the table from the next open, not the zero of the closed market after the filing. Two
filings up and one down give a mean CAR of 0.978\% that no test finds evidence in; the two during the session, a group
of their own, give 2.426\% ($t = 31.8$). Friday's filing, followed to Monday afternoon, overlaps Monday's filing: their
standardized abnormal returns are correlated over the estimation minutes they share ($\bar r = -0.008$), and the BMP t
moves from 0.600 to 0.608.

\nmsub{The files}
\fn{write\_results} writes, for each period (its name in lower case, with hyphens):
\begin{valuetable}{@{}lX@{}}
\textbf{File} & \textbf{What it holds}\\\midrule
\fn{default-events.csv} & an event a row: its CAR, standardized CAR, abnormal volume and beta\\
\fn{default-from-the-filing.csv} & a minute a row: the events known then, the AAR and CAAR with their t\\
\fn{default-from-the-next-open.csv} & the same from the next open, for the filings outside the session\\
\fn{default-groups.csv} & a group a row: its events and its tests at the end of the period, with their p-values\\
\fn{default-groups-from-the-filing.csv} & each group's AAR and CAAR a minute at a time, each row with its group\\
\fn{default-groups-from-the-next-open.csv} & the same from the next open\\
\end{valuetable}
A study of FOMC releases writes \fn{tight-moves.csv} (a release a row, an instrument a column),
\fn{tight-summary.csv} (an instrument a row: its unit, releases, mean move, standard deviation and t) and
\fn{tight-paths.csv} (the mean move a minute at a time, by instrument).

In Julia, \fn{results = analyse(read\_study(folder))} gives the same, a named tuple a period:
\begin{itemize}
  \item \fn{results[1].tests}: the tests at the end of the period, with their p-values;
  \item \fn{results[1].events}: the events' table;
  \item \fn{results[1].from\_filing} and \fn{results[1].from\_open}: the AAR and CAAR a minute at a time;
  \item \fn{results[1].curves}: each event's CAR a minute at a time;
  \item \fn{results[1].groups}: the same for each group, with its \fn{name}.
\end{itemize}
\fn{report(results)} prints them and \fn{write\_results(results, folder)} writes them.

\nmsection{The notebook}
Each study folder has a \fn{study.ipynb} that Pollis writes when the study is made, a Julia notebook in four parts:
\begin{enumerate}
  \item \textbf{What the folder holds}: \fn{study.toml}, \fn{events.csv}, \fn{prices/} and \fn{news/} (the articles
    Google News lists from a day before each filing to a day after), and what the package computes from them.
  \item \textbf{Install and run.} The cell adds PollisRegistry when it is not among the registries, adds
    PollisEventStudies when it is not installed and updates it when it is older than the study needs, then runs
    \fn{run\_study} on the notebook's own folder (or the study's path, when the notebook was moved), which prints the
    report and writes \fn{results/}.
  \item \textbf{Each event}: \fn{results[1].events}, the table of the first period's events.
  \item \textbf{Notes}: the two models of normal returns, the estimation days the BMP test needs, Yahoo Finance's 30
    days of minute prices, and the benchmark's trading only during the session.
\end{enumerate}
Events still capturing prices count up to the last complete minute, so running the notebook again as they come in
brings the results up to date. A study of FOMC releases has its own notebook, with the instruments, the contract
months, Kuttner's scaling and \fn{results[1].moves}. So does a panel event study, made on the same page from a CSV
table (Kind: Panel Event Study), which Pollis copies into the folder as \fn{panel.csv}: its notebook runs the estimators
\fn{study.toml} names (adding MixedModels.jl for the stepped wedge) and writes \fn{results/}, whose effects by event
time, pre-trends tests, Callaway and Sant'Anna's averages and stepped wedge's estimate the page then draws.

\nmsection{What the tests show}
Three synthetic filings over seven trading days, with known betas, jumps and volume; one filed after Friday's close.
\begin{valuetable}{@{}lXl@{}}
\textbf{Test} & \textbf{Built in} & \textbf{Measured}\\\midrule
Betas & 1.2, 0.8, 1.0 & 1.195, 0.828, 0.988\\
Jump at the filing & +2\%, +1\% & +1.79\%, +1.42\% the minute after\\
Jump at Monday's open & $-3$\% & $-2.97$\%\\
Thin trading, a trade every 4th minute & beta 1.0 & 0.93 (zeros in between: 0.33)\\
Abnormal volume, 5$\times$ for 60 of 120 minutes & $\log 5 / 2 = 0.805$ & 0.811 ($t = 91.5$)\\
Cross-correlation, independent noise & near 0 & $-0.008$; BMP t 0.600, KP t 0.608\\
Cross-correlation, two of three events alike & $\bar r = 1/3$, factor $\sqrt{0.4}$ & 0.33333, 0.63246\\
Generalized rank, CARs far out & large t & $t = 3.83$, $p = 0.0002$\\
Generalized rank, CARs among the rest & no evidence & $t = 0.61$, $p = 0.55$\\
FOMC: FF1 $-10$ cents on the 10th, scaled $30/20$ & 15 bp & 15 bp, to 9 decimals\\
FOMC: 5-minute bars, FF1 $-5$ cents & 5 bp from minute 5 & mean path 7.5 bp at minute 1, 10 at 5\\
\end{valuetable}
The tests also check the groups, the files written, the events still capturing and a study without events.

\nmsection{Panel event studies}
Units of a panel first treated at different times (firms adopting a rule, states passing a law, hospitals joining a
stepped-wedge trial) against units not yet or never treated. \fn{panel\_event\_study} takes any Tables.jl table, a row a
unit and a time, and gives the effect at each event time (the time since the unit was first treated) from
\fn{-leads} to \fn{lags}, against the \fn{reference} event time ($-1$ by default); event times beyond the ends count at
the ends, and units treated from the first time are left out.

\begingroup\footnotesize
\begin{verbatim}
results = panel_event_study(table; unit = "state", time = "year", outcome = "y",
    cohort = "first_treated", leads = 5, lags = 5)
report(results)
\end{verbatim}
\endgroup

\nmsub{Three estimators}
\begin{itemize}
  \item \textbf{Two-way fixed effects (TWFE).} The outcome on unit and time effects and an indicator of each event time.
    With effects that differ from cohort to cohort, its effects mix in those of other event times, before treatment
    too, with weights that can be negative.
  \item \textbf{Sun and Abraham (\nmcite{sun2021}{2021}).} Each cohort's effect at each event time against the control
    units (\fn{control = "never"}, the units never treated, or \fn{"last"}, the units treated last, before their
    treatment), averaged over the cohorts with their shares of the event time's observations.
  \item \textbf{Callaway and Sant'Anna (\nmcite{cs2021}{2021}),} \fn{callaway\_santanna}. Each cohort's effect at each
    time, ATT$(g, t)$, a difference in differences against the units never treated, not yet treated or treated last;
    from a varying or a universal base period; doubly robust (Sant'Anna and Zhao \nmcite{sz2020}{2020}), by inverse
    probability weighting or by outcome regression, with covariates; then aggregated by event time, by cohort, by
    calendar time and overall, the standard errors counting the estimation of the cohorts' shares. It is a port of
    R's did and DRDID.
\end{itemize}
Each result of \fn{panel\_event\_study} has the effects by event time, each cohort's, the pre-trends test (that the effects before the reference
event time are all zero) and the sample. \fn{weights} weighs the observations and \fn{covariates} are added to both
regressions. A study folder with \fn{source = "panel"} in \fn{study.toml} names the table, the same keywords and the estimators
(\fn{estimators}: \fn{panel\_event\_study}, \fn{callaway\_santanna} or both), the second's own keywords in a
\fn{[callaway\_santanna]} table, so that a study made on Pollis's Event Studies page can run either.

\nmsub{Standard errors: three choices}
Standard errors are clustered by \fn{cluster}, the unit by default. Three keywords set how, each with the package's
choice as the default and R's as the other:
\begin{itemize}
  \item \fn{small\_sample} (\fn{true}): Stata's small-sample factor $\frac{G}{G-1}\frac{N-1}{N-K}$, $G$ the clusters,
    $N$ the observations and $K$ the parameters, as fixest; \fn{false} leaves it out, as did.
  \item \fn{fixed\_shares} (\fn{false}): Sun and Abraham's variance adds that of the cohorts' shares, which are
    estimated too, clustered as the cohorts' effects, as did; \fn{true} holds the shares fixed,
    as fixest.
  \item \fn{distribution} (\fn{"t"}): p-values and the 95\% bands from Student's t with $G-1$ degrees of freedom, and
    the pre-trends test as $F(q, G-1)$; \fn{"normal"} takes the normal and $\chi^2(q)$, as did.
\end{itemize}

\nmsub{t or the normal}
The normal is asymptotic. The standard errors come from an influence function, which justifies the normal as the
number of clusters grows; with 50 states it is a fair approximation. The t with $G-1$ degrees of freedom is a
small-sample correction, Stata's convention, which fixest follows. With few clusters the clustered standard errors are
themselves noisy, and the t's thicker tails allow for that. With 10 to 20 clusters the difference is real: with 15, the
5\% critical value is 2.14 under $t(14)$ and 1.96 under the normal. With 50 it is 2.01, and with hundreds of clusters
the two agree.

\newpage
\nmsub{Against R}
R is the benchmark for these estimators. The tests run the package and R on the same synthetic panels: 60 units over
12 years in three cohorts and a never-treated group, with effects that differ from cohort to cohort.
\begin{valuetable}{@{}>{\raggedright\arraybackslash}p{43mm}>{\raggedright\arraybackslash}p{50mm}>{\raggedright\arraybackslash}X@{}}
\textbf{R} & \textbf{Keywords} & \textbf{Equal to 6 decimals}\\\midrule
fixest, \fn{sunab} and \fn{i()} & \fn{fixed\_shares = true} & 49 cohort and TWFE effects and standard errors, with weights, a covariate and clusters coarser than the units; 32 effects by event time and their standard errors; 14 with the last-treated control\\
did (\nmcite{cs2021}{Callaway and Sant'Anna 2021}), \fn{att\_gt} (never treated, universal base period) and \fn{aggte(type = "dynamic")}, \fn{bstrap = FALSE, cband = FALSE} & \fn{small\_sample = false, distribution = "normal"} & 33 cohort effects; 16 effects by event time with their standard errors, p-values and 95\% bands\\
did, \fn{att\_gt} and \fn{aggte} (dynamic, group, calendar, simple), \fn{bstrap = FALSE, cband = FALSE} & \fn{callaway\_santanna}, \fn{distribution = "normal"} & In six configurations (each control group, base period and method, with weights, a covariate, coarser clusters, anticipation, and with no units never treated): 164 cells, 85 effects by event time, 15 by cohort, 45 by calendar time and 24 overall, with their standard errors, to 5e-13; the pre-trends test; \fn{min\_e}, \fn{max\_e} and \fn{balance\_e}\\
\end{valuetable}
did is Sun and Abraham's estimator only with the never-treated control, a universal base period and no covariates.
Its default bands are simultaneous, from a multiplier bootstrap; \fn{bstrap = FALSE, cband = FALSE} gives the pointwise
bands compared here. Stata's eventstudyinteract takes the shares' variance observation by observation, which gives the
same variance at each event time when each unit is a cluster and no event time is beyond the ends, but no covariance
between event times. Not yet checked against R: the binned ends, the pre-trends test and the weighted last-treated
control. The bootstraps are in the next section, and datasets with gaps in the one after.

\nmsection{Bootstrap}
The bootstrap serves the panel estimators three ways: standard errors that need no formula, bands that cover all the
event times at once, and tests that hold with few clusters. Each takes \fn{draws} and \fn{seed}, to repeat the draws.
\begin{valuetable}{@{}>{\raggedright\arraybackslash}p{36mm}>{\raggedright\arraybackslash}p{44mm}>{\raggedright\arraybackslash}X@{}}
\textbf{Bootstrap} & \textbf{Keywords} & \textbf{What it gives}\\\midrule
did's multiplier bootstrap & \fn{callaway\_santanna}: \fn{bootstrap = true}; \fn{panel\_event\_study}: \fn{uniform = true, multiplier = true} & Standard errors and simultaneous bands for Callaway and Sant'Anna; did's critical values for TWFE and Sun and Abraham\\
Sup-t & \fn{panel\_event\_study}: \fn{uniform = true} & Simultaneous bands for TWFE and Sun and Abraham\\
Wild cluster bootstrap & \fn{panel\_event\_study}: \fn{wild = true}, after \fn{using WildBootTests} & p-values and 95\% intervals for each effect, and the pre-trends test, with few clusters\\
\end{valuetable}

\nmsub{Simultaneous bands}
A 95\% band at each event time covers that event time's effect 95\% of the time; drawn together on an event-study
plot, the bands cover all the effects at once far less often. Simultaneous bands cover them all together 95\% of the
time: each effect $\pm$ a critical value $c$ times its standard error, $c$ the 95\% quantile of the largest $|t|$ over
the event times. The package draws that largest $|t|$ two ways.
\begin{itemize}
  \item \textbf{Sup-t} (Montiel Olea and Plagborg-M\o{}ller \nmcite{mop2019}{2019}), the default: from the effects'
    estimated correlations, normal or Student's t as \fn{distribution} (10{,}000 draws by default).
  \item \textbf{did's multiplier bootstrap}: each cluster's influence on the effects weighed by $+1$ or $-1$ in each
    draw, each effect's draws scaled by their interquartile range over the normal's, and the pointwise value when the
    simultaneous one is smaller, as did's \fn{aggte}. For \fn{callaway\_santanna} it also gives the standard errors.
\end{itemize}
On the test panel, Sun and Abraham's 16 event times under the normal give 2.777 by sup-t, between the pointwise 1.960
and Bonferroni's 2.955. By did's multiplier bootstrap they give 2.706, and did itself 2.708 (200{,}000 draws each).
The two ways differ by the weights alone: over 60 clusters, weights of $+1$ or $-1$ give the largest $|t|$ thinner tails
than normal draws do, and normal weights on the same influences give 2.775. For \fn{callaway\_santanna}, with
100{,}000 draws, the critical values and standard errors are did's within the noise of the draws (2.743 against 2.747
for the event times, 0.1725 against 0.1729 for a cell's standard error).

\nmsub{The wild cluster bootstrap}
Clustered standard errors assume many clusters. With 10 to 30, tests on them reject too often, even with t's critical
values (Cameron, Gelbach and Miller \nmcite{cgm2008}{2008}). The wild cluster bootstrap does better: it refits the
regression with the null imposed, and in each draw multiplies each cluster's residuals by a single random weight, keeping
their correlation within the cluster; the p-value is the share of the draws whose $|t|$ is at least the data's. The
interval holds the values the test does not reject at 5\%. With \fn{wild = true}, each effect of TWFE and of Sun and
Abraham has its p-value and interval, and the pre-trends test its p-value, a Wald test with the null imposed. The weights
are $+1$ or $-1$ (\fn{wild\_weights = "rademacher"}), every one of the $2^G$ patterns of signs when there are no more
than \fn{draws}, or Webb's (\nmcite{webb2023}{2023}) six values (\fn{"webb"}), better with fewer than about 12 clusters,
whose $2^G$ patterns are then too few. Sun and Abraham's effects hold the cohorts' shares fixed. The bootstrap runs in
David Roodman's WildBootTests.jl (Roodman, MacKinnon, Nielsen and Webb \nmcite{rmnw2019}{2019}), the algorithms of
Stata's boottest, which the package loads only when asked for.

Against R's fwildclusterboot (\fn{boottest} on fixest's \fn{i()} and \fn{sunab}), on the test panel with weights, a
covariate and 12 clusters, so that all 4{,}096 patterns are drawn: the p-values are R's to 6 decimals (0.032715 for TWFE
at event time $-2$, 134 of the 4{,}096 draws; 0.966309 for Sun and Abraham). The p-value is a step function of the
value tested, and the interval's bound is where it falls below 5\%. Neither program finds that point exactly: TWFE's
upper bound at event time 3 is 1.158875, R stops its root search after 10 iterations at 1.158985, and WildBootTests
stops within a step of 5\% at 1.159517. Both are on the step past the bound, so the intervals agree to its width, here
0.0005.

\nmsection{Unbalanced datasets}
Real panels have gaps: a firm enters the sample late, leaves it, or misses a year of data. A panel is unbalanced when
some units are not observed at every time. Gaps do not bias the estimates when they are unrelated to the outcome; when
units leave because of how they fare, no estimator here corrects that.

\nmsub{TWFE and Sun and Abraham}
\fn{panel\_event\_study} needs no option. Its regression takes every observation there is, and the unit and time
effects absorb who is present when. On a panel where half the 60 units miss a year or two, it agrees with fixest to 6
decimals, effects and standard errors, with weights, a covariate and clusters coarser than the units.

\nmsub{Callaway and Sant'Anna}
Each ATT(g, t) is a change from the base period to t, which a unit missing at either time does not have. By default,
as in did, \fn{callaway\_santanna} leaves out every unit not observed at every time, and the report says how many.
With \fn{unbalanced = true}, as did's \fn{allow\_unbalanced\_panel}, it keeps them:
\begin{itemize}
  \item each cell compares the observations at t and at the base period as two cross sections, with the estimators of
    \nmcite{sz2020}{Sant'Anna and Zhao 2020} for repeated cross sections (DRDID's \fn{drdid\_rc}, \fn{std\_ipw\_did\_rc}
    and \fn{reg\_did\_rc}), each observation with its own covariates and weight;
  \item a unit's influence is the sum of its observations', so the standard errors, clustered or bootstrapped, still
    treat a unit's observations as one;
  \item the cohorts' sizes weigh each unit by the mean of its weights.
\end{itemize}
When every unit is observed at every time, the option changes nothing. Keeping the units with gaps keeps their
information, but each cell no longer differences a unit's outcome against itself, which costs precision when units
differ a lot. On the same panel the effect at event time 0, whose true value is $(0.4 + 2.0 - 1.0)/3 \approx 0.467$:
\begin{valuetable}{@{}lll@{}}
\textbf{Panel} & \textbf{Effect at 0} & \textbf{Standard error}\\\midrule
The 30 units with gaps left out (default) & 0.380 & 0.259\\
All 60 units, \fn{unbalanced = true} & 0.494 & 0.196\\
The full panel, without gaps & 0.468 & 0.187\\
\end{valuetable}
Against did 2.5.1 with \fn{allow\_unbalanced\_panel = TRUE}, in five configurations (each method, the not-yet-treated
control, a universal base period, weights, a covariate and coarser clusters), the cells, the effects by event time, by
cohort and overall, their standard errors and the pre-trends test are equal to 6 decimals.

\nmsection{Stepped wedges}
In a stepped-wedge trial every cluster starts the treatment, each at a time drawn at random. \fn{stepped\_wedge}, after
\fn{using MixedModels}, fits the mixed model of \nmcite{hh2007}{Hussey and Hughes (2007)}: a fixed effect for each time,
a random intercept for each cluster, and one for each cluster at each time, so that outcomes in a cluster are more alike
at the same time than at different times. A treated observation's exposure time is the time since its cluster's first
treatment, plus 1.

\begingroup\footnotesize
\begin{verbatim}
using MixedModels
result = stepped_wedge(table; cluster = "clinic", time = "period", treatment = "trt",
    outcome = "y", effect = "ETI")
\end{verbatim}
\endgroup

\begin{itemize}
  \item \fn{effect = "IT"}, immediate treatment: one effect at every exposure time, the model of Hussey and Hughes;
  \item \fn{"ETI"}, exposure time indicators: an effect at each exposure time (\nmcite{kvxhh2022}{Kenny et al. 2022});
  \item \fn{"TEH"}, treatment effect heterogeneity: a random effect at each exposure time about a mean effect.
\end{itemize}
The estimand is the TATE, the effects averaged over a range of exposure times (all by default), or a PTE, the effect at
one exposure time. Gaussian outcomes are fitted by REML, binary ones, or successes out of trials, by the Laplace
approximation; the result also has the effects by exposure time, the standard deviations of the random effects and,
for a gaussian outcome, the intracluster correlation. The model runs in MixedModels.jl, which the package loads only when
asked for. In a study folder, \fn{estimators = ["stepped\_wedge"]} in \fn{study.toml}, with the keywords in a
\fn{[stepped\_wedge]} table (a TATE's exposure times as \fn{exposure = [2, 5]}), and \fn{run\_study} writes
\fn{sw-effects.csv} and \fn{sw-estimate.csv}.

When the effect grows with exposure time, IT's single effect mixes the early, small effects with the later ones and can
be far from their average. In the test simulation, 18 clusters over 7 times whose effect rises from 0.2 at exposure
time 1 to 1.0 at 6, an average of 0.683:
\begin{valuetable}{@{}lll@{}}
\textbf{Effect} & \textbf{TATE} & \textbf{Standard error}\\\midrule
IT & 0.181 & 0.120\\
ETI & 0.693 & 0.173\\
TEH & 0.524 & 0.185\\
\end{valuetable}

\nmsub{Against R}
The code is a port of the mixed-model part of R's steppedwedge 1.0.0 (\fn{analyze(method = "mixed")}, calendar time
categorical). On steppedwedge's simulated example (15 clusters, 5 sequences, 6 times) and on the simulation above, the
gaussian IT and ETI estimates and standard errors are R's within $10^{-6}$. The binomial fits and TEH are within
$5\times10^{-5}$: a variance near zero, as that of the clusters at a time often is, leaves the likelihood flat there,
and R stops sooner, at a likelihood no higher than here ($-1376.65973195$ against $-1376.65973191$). Under TEH, the
standard error of a PTE takes the conditional variance of the exposure time asked for; steppedwedge takes that of the
exposure time before it (0.220108 against R's 0.213374 at exposure time 4).

\nmsection{Case studies}
\nmsub{A synthetic panel: TWFE's bias}
The panel of the tests: 60 units over the years 1 to 12, 15 first treated in year 4, 15 in 7, 15 in 9 and 15 never,
with effects that differ from cohort to cohort: $0.4(l+1)$ at event time $l$ for the first cohort, 2 for the second and
$-1$ for the third, and none before treatment. The true effect at an event time is the average over the cohorts seen
there. With every event time kept (\fn{leads = 8, lags = 8}), so that none is binned:
\begingroup\footnotesize
\begin{verbatim}
results = panel_event_study(panel; unit = "unit", time = "year", outcome = "y",
    cohort = "first", leads = 8, lags = 8)
cs = callaway_santanna(panel; unit = "unit", time = "year", outcome = "y", cohort = "first")
\end{verbatim}
\endgroup
\begin{valuetable}{@{}rrrrr@{}}
\textbf{Event time} & \textbf{True} & \textbf{TWFE} & \textbf{Sun and Abraham} & \textbf{Callaway and Sant'Anna}\\\midrule
$-8$ & 0 & 1.318 & $-0.002$ & \\
$-6$ & 0 & 0.550 & $-0.015$ & $-0.012$\\
$-4$ & 0 & 0.295 & $-0.015$ & $-0.001$\\
$-3$ & 0 & 0.331 & 0.005 & 0.003\\
$-2$ & 0 & 0.044 & 0.004 & 0.000\\
0 & 0.467 & 0.458 & 0.468 & 0.468\\
2 & 0.733 & 1.019 & 0.741 & 0.741\\
4 & 2.000 & 1.774 & 2.015 & 2.015\\
6 & 2.800 & 2.815 & 2.847 & 2.847\\
8 & 3.600 & 4.167 & 3.618 & 3.618\\
\end{valuetable}
TWFE finds pre-trends where there are none (its test: $F(7, 59) = 26.5$, $p < 10^{-15}$) and effects up to 0.57 from
the true ones after treatment: each of its event times takes in the effects of the others, from other cohorts. Sun and
Abraham's and Callaway and Sant'Anna's are within the noise of the truth, and Sun and Abraham's pre-trends test does
not reject ($p = 0.99998$). After treatment the two are the same numbers, each cohort against the units never treated; before
it, Callaway and Sant'Anna compares each year with the one before (the varying base period), Sun and Abraham with
event time $-1$.

\nmsub{The minimum wage: a published table}
\fn{mpdta}, in PollisDatasets, is the did package's sample of 500 counties from 2003 to 2007, of which some states
raised the minimum wage in 2004, 2006 or 2007; the outcome is the log of teen employment (\fn{lemp}). Callaway and
Sant'Anna's own vignette (did 2.5.1, \emph{Getting Started with the did Package}) publishes the effects without
covariates, against the units never treated, doubly robust:
\begingroup\footnotesize
\begin{verbatim}
using PollisDatasets
mpdta = dataset("mpdta")
cs = callaway_santanna(mpdta; unit = "countyreal", time = "year", outcome = "lemp",
    cohort = "first_treat", distribution = "normal")
\end{verbatim}
\endgroup
\begin{valuetable}{@{}lrrr@{}}
\textbf{Effect} & \textbf{did, published} & \textbf{Package} & \textbf{Package's standard error}\\\midrule
Event time $-3$ & 0.0305 & 0.0305 & 0.0150\\
Event time $-2$ & $-0.0006$ & $-0.0006$ & 0.0133\\
Event time $-1$ & $-0.0245$ & $-0.0245$ & 0.0142\\
Event time 0 & $-0.0199$ & $-0.0199$ & 0.0118\\
Event time 1 & $-0.0510$ & $-0.0510$ & 0.0169\\
Event time 2 & $-0.1373$ & $-0.1373$ & 0.0364\\
Event time 3 & $-0.1008$ & $-0.1008$ & 0.0344\\
Overall, by event time & $-0.0772$ & $-0.0772$ & 0.0200\\
Overall, \fn{balance = 1} & $-0.0288$ & $-0.0288$ & 0.0137\\
Pre-trends test, p-value & 0.16812 & 0.16812 & \\
\end{valuetable}
The 12 cohort-by-year effects are the published ones to the same 4 decimals, as are the four effects with
\fn{balance = 1}. The published standard errors are did's default, a multiplier bootstrap, whose draws differ from run
to run (0.0203 for the overall effect by event time); the package's are the analytic ones, those of did with
\fn{bstrap = FALSE}. TWFE and Sun and Abraham on the same data (\fn{leads = 3, lags = 3}) agree on the sign and the
growth of the effect: $-0.0188$ and $-0.0209$ at event time 0, $-0.1329$ and $-0.1418$ at 2.

\nmsub{No-fault divorce: Goodman-Bacon's example}
\fn{divorce}, in PollisDatasets, is the age-standardized female suicide rate (per million women) in 49 states from 1964
to 1996, and whether each had no-fault (unilateral) divorce (\nmcite{sw2006}{Stevenson and Wolfers 2006}), as Thomas
Goldring's ddtiming (MIT) ships it. It is Goodman-Bacon's (\nmcite{gb2021}{2021}) example, not Sun and Abraham's. 12
cohorts adopted the law between 1969 and 1985, 8 states had it before 1964 and 5 had not by 1996. The regression of
the rate on the indicator with state and year effects gives the published $-3.08$ ($-3.0799$, in R's \fn{lm}), and
Goodman-Bacon's decomposition of it into its two-by-two comparisons (R's bacondecomp) shows where it comes from:
\begin{valuetable}{@{}lrr@{}}
\textbf{Comparison} & \textbf{Weight} & \textbf{Average estimate}\\\midrule
Later treated against the states treated before 1964 & 0.384 & $-7.04$\\
Later treated against earlier treated & 0.265 & 3.51\\
Treated against never treated & 0.240 & $-5.33$\\
Earlier treated against later treated & 0.111 & $-0.19$\\
\end{valuetable}
A quarter of the weight compares states with states already treated, whose effect is still changing, and has the
wrong sign. With his controls (income, homicides, welfare) the estimate is $-2.52$, from a term with 0.5\% of the
weight; the comparisons across cohorts barely move. The package's estimators leave out the 8 states treated from the
start, which have no year before treatment, and with them the comparisons that carry 38\% of TWFE's weight:
\begingroup\footnotesize
\begin{verbatim}
divorce = dataset("divorce")
results = panel_event_study(divorce; unit = "state", time = "year", outcome = "asmrs",
    treatment = "treat", leads = 5, lags = 10)
cs = callaway_santanna(divorce; unit = "state", time = "year", outcome = "asmrs",
    treatment = "treat", leads = 5, lags = 10)
\end{verbatim}
\endgroup
\begin{valuetable}{@{}rrrr@{}}
\textbf{Event time} & \textbf{TWFE} & \textbf{Sun and Abraham} & \textbf{Callaway and Sant'Anna}\\\midrule
$-5$ & $-2.70$ (3.18) & $-0.76$ (3.66) & $-2.11$ (3.08)\\
$-3$ & $-2.46$ (2.85) & $-0.94$ (3.17) & $-0.89$ (2.24)\\
$-2$ & $-0.36$ (2.38) & $-0.02$ (2.85) & 1.51 (2.66)\\
0 & 0.56 (2.64) & 1.61 (3.06) & $-0.76$ (2.76)\\
2 & $-0.62$ (3.97) & 0.85 (4.38) & $-3.59$ (4.21)\\
4 & $-1.60$ (2.89) & 0.00 (3.22) & $-4.88$ (2.89)\\
6 & $-1.94$ (3.76) & $-1.81$ (4.10) & $-6.30$ (3.83)\\
8 & $-5.79$ (3.23) & $-6.36$ (3.95) & $-9.95$ (3.59)\\
10 & $-8.87$ (3.72) & $-9.35$ (4.40) & $-11.18$ (3.84)\\
\end{valuetable}
Standard errors in brackets; event time 10 takes in the later ones, and $-5$ the earlier. No estimator finds
pre-trends (TWFE's test $p = 0.58$, Sun and Abraham's 0.89), and all three find a decline in female suicide that
grows with the years since the law, small in the first years, about 6 to 11 deaths per million women ten years on:
the single number of TWFE, $-3.08$, averages a growing effect with comparisons of the wrong sign. Without
\fn{leads} and \fn{lags}, Callaway and Sant'Anna's average over all the cohorts and years after treatment
(\fn{cs.overall}, simple) is $-10.22$ (3.09), did's to the digits shown, against
the 5 states never treated, the only clean controls, which is why its standard errors stay large.

\newpage
\nmsection{Still open}
\begin{itemize}
  \item \textbf{A published intraday table.} The panel estimators match a published table (the minimum wage, above);
    the intraday tests are all synthetic. The package needs an intraday study whose events and
    results are published and whose prices can be had today, to be tested against as the other packages are.
  \item \textbf{Estimation days.} Pollis captures them for the market model only, so market-adjusted studies have no
    standardized tests (BMP, its adjustment and the generalized rank test) unless the prices before the period happen
    to be there.
\end{itemize}

\nmsection{References}
{\footnotesize
\nmref{bm2003}{Bertrand, M. and Mullainathan, S. (2003). Enjoying the quiet life? Corporate governance and managerial preferences. \emph{Journal of Political Economy} 111(5), 1043--1075. \nmdoi{10.1086/376950}}
\nmref{bmp1991}{Boehmer, E., Musumeci, J. and Poulsen, A.\,B. (1991). Event-study methodology under conditions of event-induced variance. \emph{Journal of Financial Economics} 30(2), 253--272. \nmdoi{10.1016/0304-405X(91)90032-F}}
\nmref{blm2022}{Braghieri, L., Levy, R. and Makarin, A. (2022). Social media and mental health. \emph{American Economic Review} 112(11), 3660--3693. \nmdoi{10.1257/aer.20211218}}
\nmref{cs2021}{Callaway, B. and Sant'Anna, P.\,H.\,C. (2021). Difference-in-differences with multiple time periods. \emph{Journal of Econometrics} 225(2), 200--230.\newline \nmdoi{10.1016/j.jeconom.2020.12.001}}
\nmref{cgm2008}{Cameron, A.\,C., Gelbach, J.\,B. and Miller, D.\,L. (2008). Bootstrap-based improvements for inference with clustered errors. \emph{Review of Economics and Statistics} 90(3), 414--427. \nmdoi{10.1162/rest.90.3.414}}
\nmref{ffjr1969}{Fama, E.\,F., Fisher, L., Jensen, M.\,C. and Roll, R. (1969). The adjustment of stock prices to new information. \emph{International Economic Review} 10(1), 1--21.}
\nmref{gb2021}{Goodman-Bacon, A. (2021). Difference-in-differences with variation in treatment timing. \emph{Journal of Econometrics} 225(2), 254--277.\newline \nmdoi{10.1016/j.jeconom.2021.03.014}}
\nmref{hh2007}{Hussey, M.\,A. and Hughes, J.\,P. (2007). Design and analysis of stepped wedge cluster randomized trials. \emph{Contemporary Clinical Trials} 28, 182--191. \nmdoi{10.1016/j.cct.2006.05.007}}
\nmref{jjp2016}{Jackson, C.\,K., Johnson, R.\,C. and Persico, C. (2016). The effects of school spending on educational and economic outcomes: evidence from school finance reforms. \emph{Quarterly Journal of Economics} 131(1), 157--218. \nmdoi{10.1093/qje/qjv036}}
\nmref{kvxhh2022}{Kenny, A., Voldal, E.\,C., Xia, F., Heagerty, P.\,J. and Hughes, J.\,P. (2022). Analysis of stepped wedge cluster randomized trials in the presence of a time-varying treatment effect. \emph{Statistics in Medicine} 41(22), 4311--4339. \nmdoi{10.1002/sim.9511}}
\nmref{kp2010}{Kolari, J.\,W. and Pynn\"onen, S. (2010). Event study testing with cross-sectional correlation of abnormal returns. \emph{Review of Financial Studies} 23(11), 3996--4025. \nmdoi{10.1093/rfs/hhq072}}
\nmref{kp2011}{Kolari, J.\,W. and Pynn\"onen, S. (2011). Nonparametric rank tests for event studies. \emph{Journal of Empirical Finance} 18(5), 953--971.\newline \nmdoi{10.1016/j.jempfin.2011.08.003}}
\nmref{kuttner2001}{Kuttner, K.\,N. (2001). Monetary policy surprises and interest rates: evidence from the Fed funds futures market. \emph{Journal of Monetary Economics} 47(3), 523--544. \nmdoi{10.1016/S0304-3932(01)00055-1}}
\nmref{lcg2017}{Lopez Bernal, J., Cummins, S. and Gasparrini, A. (2017). Interrupted time series regression for the evaluation of public health interventions: a tutorial. \emph{International Journal of Epidemiology} 46(1), 348--355. \nmdoi{10.1093/ije/dyw098}}
\nmref{mackinlay1997}{MacKinlay, A.\,C. (1997). Event studies in economics and finance. \emph{Journal of Economic Literature} 35(1), 13--39.}
\nmref{mjw2021}{Miller, S., Johnson, N. and Wherry, L.\,R. (2021). Medicaid and mortality: new evidence from linked survey and administrative data. \emph{Quarterly Journal of Economics} 136(3), 1783--1829. \nmdoi{10.1093/qje/qjab004}}
\nmref{mop2019}{Montiel Olea, J.\,L. and Plagborg-M\o{}ller, M. (2019). Simultaneous confidence bands: Theory, implementation, and an application to SVARs. \emph{Journal of Applied Econometrics} 34(1), 1--17. \nmdoi{10.1002/jae.2656}}
\nmref{rmnw2019}{Roodman, D., MacKinnon, J.\,G., Nielsen, M.\,\O{}. and Webb, M.\,D. (2019). Fast and wild: Bootstrap inference in Stata using boottest. \emph{Stata Journal} 19(1), 4--60. \nmdoi{10.1177/1536867X19830877}}
\nmref{rs2023}{Roth, J. and Sant'Anna, P.\,H.\,C. (2023). Efficient estimation for staggered rollout designs. \emph{Journal of Political Economy Microeconomics} 1(4), 669--709. \nmdoi{10.1086/726581}}
\nmref{sz2020}{Sant'Anna, P.\,H.\,C. and Zhao, J. (2020). Doubly robust difference-in-differences estimators. \emph{Journal of Econometrics} 219(1), 101--122.\newline \nmdoi{10.1016/j.jeconom.2020.06.003}}
\nmref{sw2006}{Stevenson, B. and Wolfers, J. (2006). Bargaining in the shadow of the law: Divorce laws and family distress. \emph{Quarterly Journal of Economics} 121(1), 267--288.\newline \nmdoi{10.1162/qjec.2006.121.1.267}}
\nmref{sun2021}{Sun, L. and Abraham, S. (2021). Estimating dynamic treatment effects in event studies with heterogeneous treatment effects. \emph{Journal of Econometrics} 225(2), 175--199. \nmdoi{10.1016/j.jeconom.2020.09.006}}
\nmref{webb2023}{Webb, M.\,D. (2023). Reworking wild bootstrap-based inference for clustered errors. \emph{Canadian Journal of Economics} 56(3), 839--858. \nmdoi{10.1111/caje.12661}}
\nmref{software}{Software: Kaspereit, T., eventstudy2 (Stata); EventStudy.jl (MIT), \url{https://github.com/florian-st/EventStudy.jl}; Berg\'e, L., fixest (R); Callaway, B. and Sant'Anna, P.\,H.\,C., did (R); Sant'Anna, P.\,H.\,C. and Zhao, J., DRDID (R); Sun, L., eventstudyinteract (Stata); Roodman, D., WildBootTests.jl (MIT), \url{https://github.com/droodman/WildBootTests.jl}; Fischer, A. and Roodman, D., fwildclusterboot (R).}
}

\nmcredit
\end{document}
